\documentclass[10pt,a4paper]{amsart}
\usepackage{amsmath,amssymb,amsthm,graphicx}
% Retained minimal preamble for the Stacks Project chapter more-algebra.tex (More on Algebra).
% The chapter file pinned-original-source.tex carries only \input{preamble} on its first line, so no
% preamble text lives inside the chapter itself.  The Stacks Project uses one single shared preamble for
% every chapter; the snapshot stages it identically as sources/stacks-chow/preamble.tex and
% sources/stacks-etale-cohomology/preamble.tex.  The block below is copied mechanically, line for line,
% from sources/stacks-chow/preamble.tex lines 14-18 (the Xy-pic setup) and 154-217 (the native theorem
% declarations together with the mathematical macros).  Nothing was added, removed or redefined; the
% wrapper class and the amsmath/amssymb/amsthm/graphicx base packages are supplied by the assembler.
% Chapter: more-algebra.tex -- More on Algebra.
% For commutative diagrams we use Xy-pic
\usepackage[all]{xy}

% We use 2cell for 2-commutative diagrams.
\xyoption{2cell}

% Theorem environments.
%
\theoremstyle{plain}
\newtheorem{theorem}[subsection]{Theorem}
\newtheorem{proposition}[subsection]{Proposition}
\newtheorem{lemma}[subsection]{Lemma}

\theoremstyle{definition}
\newtheorem{definition}[subsection]{Definition}
\newtheorem{example}[subsection]{Example}
\newtheorem{exercise}[subsection]{Exercise}
\newtheorem{situation}[subsection]{Situation}

\theoremstyle{remark}
\newtheorem{remark}[subsection]{Remark}
\newtheorem{remarks}[subsection]{Remarks}

\numberwithin{equation}{subsection}

% Macros
%
\def\lim{\mathop{\mathrm{lim}}\nolimits}
\def\colim{\mathop{\mathrm{colim}}\nolimits}
\def\Spec{\mathop{\mathrm{Spec}}}
\def\Hom{\mathop{\mathrm{Hom}}\nolimits}
\def\Ext{\mathop{\mathrm{Ext}}\nolimits}
\def\SheafHom{\mathop{\mathcal{H}\!\mathit{om}}\nolimits}
\def\SheafExt{\mathop{\mathcal{E}\!\mathit{xt}}\nolimits}
\def\Sch{\mathit{Sch}}
\def\Mor{\mathop{\mathrm{Mor}}\nolimits}
\def\Ob{\mathop{\mathrm{Ob}}\nolimits}
\def\Sh{\mathop{\mathit{Sh}}\nolimits}
\def\NL{\mathop{N\!L}\nolimits}
\def\CH{\mathop{\mathrm{CH}}\nolimits}
\def\proetale{{pro\text{-}\acute{e}tale}}
\def\etale{{\acute{e}tale}}
\def\QCoh{\mathit{QCoh}}
\def\Ker{\mathop{\mathrm{Ker}}}
\def\Im{\mathop{\mathrm{Im}}}
\def\Coker{\mathop{\mathrm{Coker}}}
\def\Coim{\mathop{\mathrm{Coim}}}

% Boxtimes
%
\DeclareMathSymbol{\boxtimes}{\mathbin}{AMSa}{"02}

%
% Macros for moduli stacks/spaces
%
\def\QCohstack{\mathcal{QC}\!\mathit{oh}}
\def\Cohstack{\mathcal{C}\!\mathit{oh}}
\def\Spacesstack{\mathcal{S}\!\mathit{paces}}
\def\Quotfunctor{\mathrm{Quot}}
\def\Hilbfunctor{\mathrm{Hilb}}
\def\Curvesstack{\mathcal{C}\!\mathit{urves}}
\def\Polarizedstack{\mathcal{P}\!\mathit{olarized}}
\def\Complexesstack{\mathcal{C}\!\mathit{omplexes}}
% \Pic is the operator that assigns to X its picard group, usage \Pic(X)
% \Picardstack_{X/B} denotes the Picard stack of X over B
% \Picardfunctor_{X/B} denotes the Picard functor of X over B
\def\Pic{\mathop{\mathrm{Pic}}\nolimits}
\def\Picardstack{\mathcal{P}\!\mathit{ic}}
\def\Picardfunctor{\mathrm{Pic}}
\def\Deformationcategory{\mathcal{D}\!\mathit{ef}}
\usepackage{cleveref}
\crefname{originallemma}{original lemma}{original lemmas}
\Crefname{originallemma}{Original lemma}{Original lemmas}
\makeatletter
\newlabel{more-algebra-lemma-join-quasi-regular-H1-regular}{{0667}{source}{}{sourcepointer-1}{}}
\newlabel{more-algebra-lemma-join-quasi-regular-H1-regular@cref}{{[originallemma][0][]0667}{[1][0][]source}{}{}{}}
\makeatother
\crefname{originallemma}{original lemma}{original lemmas}
\Crefname{originallemma}{Original lemma}{Original lemmas}
\makeatletter
\newlabel{lemma-H1-regular-quasi-regular}{{062I}{source}{}{sourcepointer-2}{}}
\newlabel{lemma-H1-regular-quasi-regular@cref}{{[originallemma][0][]062I}{[1][0][]source}{}{}{}}
\makeatother
\crefname{originallemma}{original lemma}{original lemmas}
\Crefname{originallemma}{Original lemma}{Original lemmas}
\makeatletter
\newlabel{lemma-H1-regular-in-quotient}{{0665}{source}{}{sourcepointer-3}{}}
\newlabel{lemma-H1-regular-in-quotient@cref}{{[originallemma][0][]0665}{[1][0][]source}{}{}{}}
\makeatother
\title{A quasi-regular sequence followed by a sequence that is H1-regular modulo the ideal it generates is quasi-regular}
\author{The Stacks Project Authors}
\date{stacks-project:more-algebra.tex:local-snapshot-20261012; GNU FDL 1.2 or later}
\begin{document}
\maketitle
\setcounter{subsection}{1}
\noindent\textit{Editorial note.} All delivered mathematical text is the author's own, copied from the pinned Stacks Project chapter file more-algebra.tex; no statement or proof step was rewritten, reordered or removed.\par\medskip
\noindent\textit{Editorial note.} The chapter file pinned-original-source.tex carries only \textbackslash{}input\{preamble\} on its first line, so no preamble text lives inside the chapter. The retained file sources/stacks-more-algebra/retained-preamble.tex is a mechanical, line-for-line copy of the project's single shared preamble as staged in the snapshot, namely sources/stacks-chow/preamble.tex lines 14-18 (Xy-pic setup) and 154-217 (theorem declarations and mathematical macros). Nothing was added, removed or redefined, no package was added, and the excerpt serves only as preamble, never as a span coordinate.\par\medskip
\noindent\textit{Editorial note.} Both counted spans start at the \textbackslash{}begin line of their environment, so the statement and the author's proof are delivered whole; surrounding prose of the section is deliberately outside the counted spans and is carried as standard prerequisites instead.\par\medskip
\noindent\textit{Editorial note.} The citation fields (source\_id, source\_version, authors, source\_url, license.name, source\_title) are identical to those of the cards already published from this chapter, so that the batch selector treats them as one source instead of deferring them as a same-source clash.\par\medskip
\noindent\textit{Allowed standard prerequisites (editorial).}\begin{itemize}
\item Quasi-regularity of a sequence and H\_1-regularity of a sequence, the two regularity notions of the section. defined in the discussion of the same section, source lines 7100-7298, before the counted statement. Supplies the hypothesis on f\_1 to f\_n, the hypothesis on the images of g\_1 to g\_m, and the goal for the juxtaposed sequence.
\item For a quasi-regular sequence, consecutive quotients of powers of the ideal are direct sums of copies of the quotient by the ideal. used at proof lines 7540-7542 where I\textasciicircum{}\{d-1\}/I\textasciicircum{}d is identified as a direct sum of copies of A/I. Makes the first complex in the displayed short exact sequence a direct sum of Koszul complexes over A/I, hence acyclic in degree one.
\item The intersection lemma relating an ideal, a sequence and its powers. the author cites this at proof lines 7566-7569 as Lemma lemma-H1-regular-in-quotient (tag 0665), the preceding key lemma of the same section. Converts the membership obtained after applying quasi-regularity of the g's into the membership in I\textasciicircum{}\{d-e\}(g\_1, ..., g\_m) needed for the final step.
\end{itemize}\medskip
\ifdefined\hypertarget\hypertarget{sourcepointer-1}{}\fi
\noindent\textit{Source reference.} Original lemma 0667: focused-seeds/AGP/sources/stacks-snapshot-20261011/tags.txt, source lines 7926--7926. The lemma's own canonical tag, quoted so that readers can resolve the label inside the delivered text; it points at the counted statement itself and carries no proof content.\par\medskip
\ifdefined\hypertarget\hypertarget{sourcepointer-2}{}\fi
\noindent\textit{Source reference.} Original lemma 062I: focused-seeds/AGP/sources/stacks-snapshot-20261011/tags.txt, source lines 7797--7797. Navigation pointer to the chapter's lemma identifying H\_1-regularity with quasi-regularity, which the author uses to pass between the two notions; the inductions over the powers I\textasciicircum{}d remain inside the counted proof span.\par\medskip
\ifdefined\hypertarget\hypertarget{sourcepointer-3}{}\fi
\noindent\textit{Source reference.} Original lemma 0665: focused-seeds/AGP/sources/stacks-snapshot-20261011/tags.txt, source lines 7924--7924. Navigation pointer to the immediately preceding intersection lemma of the same chapter, invoked at the last step of the proof; that statement is counted by the sibling card S1-algebra-H1-regular-in-quotient of this lane and nothing here re-proves it.\par\medskip
\begin{lemma}
\label{lemma-join-quasi-regular-H1-regular}
Let $A$ be a ring. Let $I$ be an ideal generated by a quasi-regular
sequence $f_1, \ldots, f_n$ in $A$. Let $g_1, \ldots, g_m \in A$ be
elements whose images $\overline{g}_1, \ldots, \overline{g}_m$ form an
$H_1$-regular sequence in $A/I$. Then $f_1, \ldots, f_n, g_1, \ldots, g_m$
is a quasi-regular sequence in $A$.
\end{lemma}

\begin{proof}
We claim that $g_1, \ldots, g_m$ forms an $H_1$-regular sequence in
$A/I^d$ for every $d$. By induction assume that this holds in
$A/I^{d - 1}$. We have a short exact sequence of complexes
$$
0 \to K_\bullet(A, g_\bullet) \otimes_A I^{d - 1}/I^d
\to K_\bullet(A/I^d, g_\bullet) \to
K_\bullet(A/I^{d - 1}, g_\bullet) \to 0
$$
Since $f_1, \ldots, f_n$ is quasi-regular we see that the first complex
is a direct sum of copies of $K_\bullet(A/I, g_1, \ldots, g_m)$
hence acyclic in degree $1$. By induction hypothesis the last complex is
acyclic in degree $1$. Hence also the middle complex is.
In particular, the sequence $g_1, \ldots, g_m$ forms a quasi-regular
sequence in $A/I^d$ for every $d \geq 1$, see
Lemma \ref{lemma-H1-regular-quasi-regular}.
Now we are ready to prove that $f_1, \ldots, f_n, g_1, \ldots, g_m$
is a quasi-regular sequence in $A$.
Namely, set $J = (f_1, \ldots, f_n, g_1, \ldots, g_m)$ and suppose
that (with multinomial notation)
$$
\sum\nolimits_{|N| + |M| = d} a_{N, M} f^N g^M \in J^{d + 1}
$$
for some $a_{N, M} \in A$. We have to show that $a_{N, M} \in J$
for all $N, M$. Let $e \in \{0, 1, \ldots, d\}$. Then
$$
\sum\nolimits_{|N| = d - e, \ |M| = e} a_{N, M} f^N g^M \in
(g_1, \ldots, g_m)^{e + 1} + I^{d - e + 1}
$$
Because $g_1, \ldots, g_m$ is a quasi-regular sequence in $A/I^{d - e + 1}$
we deduce
$$
\sum\nolimits_{|N| = d - e} a_{N, M} f^N \in
(g_1, \ldots, g_m) + I^{d - e + 1}
$$
for each $M$ with $|M| = e$. By
Lemma \ref{lemma-H1-regular-in-quotient}
applied to $I^{d - e}/I^{d - e + 1}$ in the ring $A/I^{d - e + 1}$
this implies $\sum_{|N| = d - e} a_{N, M} f^N \in I^{d - e}(g_1, \ldots, g_m)$.
Since $f_1, \ldots, f_n$ is quasi-regular in $A$ this implies
that $a_{N, M} \in J$ for each $N, M$ with $|N| = d - e$ and $|M| = e$.
This proves the lemma.
\end{proof}

\end{document}
